\documentclass[10pt]{letter}
\usepackage[a4paper, margin=0.85in]{geometry}
\usepackage[T1]{fontenc}
\setlength{\parskip}{5pt}
\signature{Harsh Bordekar \\ harshbordekar@gmail.com}
\begin{document}
\begin{letter}{Airbus Operations GmbH \\ Hamburg}
\opening{Dear Hiring Team,}

I am writing to apply for the F\&DT Lead Engineer for Fuselage Primary Structure position in Hamburg/Bremen. I am a structural analysis and simulation engineer with over seven years in aerospace structural integrity, spanning fatigue and damage tolerance of metallic and composite structures, and my work centers on exactly what this role asks for: translating damage models into allowable-limit philosophies and the stress justification that structural approval depends on.

During my PhD at DLR and TU Braunschweig I derived a certification-aligned failure criterion, a leakage envelope, that maps combined strain and load states to failure onset. This defines an allowable-damage-limit philosophy and provides the failure-criterion definition and stress justification required for structural approval of composite pressure structures. Earlier, as a Research Intern in Metallic Fatigue and Damage Tolerance at DLR, I derived allowable defect sizes for a target fatigue life using small-defect fatigue models (Murakami, El-Haddad) and translated fracture-mechanics and fatigue criteria into allowable-defect acceptance limits for manufactured metallic parts. Both are the same kind of damage-tolerance-philosophy work your posting describes, applied to different structures.

Beyond deriving these criteria, I regularly provide sign-off-style technical review of externally and student-produced structural analysis work and define the technical scope for that work. I have instructed and supervised 60+ Master's students in FEM theory and ABAQUS-based modelling. I work daily with the FE solvers (Nastran/Patran, ABAQUS) and analytical stress methods (Niu, Bruhn) that sit alongside F\&DT sizing, and I build Python-based automation for analysis and reporting. I am also genuinely interested in applying AI to improve engineering workflows.

To be transparent about fit: my leadership experience is technical rather than formal line management, I have not managed direct reports in a company structure, and my evidence of leading and reviewing others' work comes from instructing and supervising university students and from sign-off review of externally produced analyses. I have not developed or validated in-service repair schemes, and I have not used Hyperworks; my FE tool experience is Nastran/Patran and ABAQUS. My German is currently at B1 level. I would rather name these gaps now than have them surface later, and I am confident the underlying F\&DT technical depth and review discipline transfer quickly to a fuselage primary structure context.

I am motivated by the chance to bring rigorous, certification-oriented F\&DT thinking to Airbus's fuselage primary structure, and I would welcome the opportunity to discuss how my background in damage tolerance philosophy, model verification, and technical review can contribute to your team.

Thank you for your consideration. I look forward to discussing how I can contribute.

\closing{Sincerely,}
\end{letter}
\end{document}
